% arXiv-ready preprint. Compile: pdflatex khinchin-2m.tex (twice).
% Suggested categories: math.NT (primary), cross-list math.NA / cs.MS.
\documentclass[11pt]{amsart}
\usepackage[T1]{fontenc}
\usepackage[utf8]{inputenc}
\usepackage{amsmath,amssymb}
\usepackage{booktabs}
\usepackage[hidelinks]{hyperref}
\usepackage{microtype}

\newcommand{\Kzero}{K_0}

\title{Khinchin's constant to two million decimal digits}
\author{Reza K Ghazi}
\email{reza@brightarcadia.com}
\urladdr{https://orcid.org/0009-0009-3286-4675}
\date{August 27, 2026}

\subjclass[2020]{Primary 11Y60; Secondary 11K50, 65G30}
\keywords{Khinchin's constant, continued fractions, arbitrary-precision
arithmetic, interval arithmetic, ball arithmetic, FLINT, Arb}

\begin{document}

\begin{abstract}
We report a computation of Khinchin's constant
$\Kzero = 2.68545200106530644530\ldots$ to $2{,}000{,}000$ decimal
digits, doubling the largest computation previously on record
($10^6$ digits, Sim\'o, 2016). The computation evaluates an
accelerated form of the Bailey--Borwein--Crandall zeta series in
FLINT/Arb ball arithmetic, so every printed digit is certified
correctly rounded rather than heuristically estimated. It completed in
8\,h\,47\,m on a single 24-thread desktop machine with a peak of
55.8\,GiB of memory; the first $10^6$ digits agree digit for digit
with Sim\'o's independent computation. The software, the digits, and
the raw run telemetry are open source and permanently archived.
\end{abstract}

\maketitle

\section{Introduction}

Write a real number $x$ as a simple continued fraction
\[
x = a_0 + \cfrac{1}{a_1 + \cfrac{1}{a_2 + \cfrac{1}{a_3 + \dotsb}}},
\qquad a_i \in \mathbb{N}.
\]
Khinchin \cite{khinchin1935} proved that for almost every real number
(in the sense of Lebesgue measure) the geometric mean of the partial
quotients converges to a universal constant independent of $x$:
\[
\lim_{n\to\infty} \bigl(a_1 a_2 \cdots a_n\bigr)^{1/n}
= \Kzero
= \prod_{k=1}^{\infty} \Bigl(1 + \frac{1}{k(k+2)}\Bigr)^{\log_2 k}
= 2.685452001\ldots
\]
The constant is entry A002210 of the OEIS \cite{oeis}. Despite its
central place in the metric theory of continued fractions, almost
nothing is known about $\Kzero$ as a number: it is not known whether it
is irrational, and it is not known whether its own continued fraction
expansion satisfies Khinchin's theorem (numerically it appears to).

Unlike $\pi$, for which hypergeometric series with rational terms admit
binary splitting and hence quasi-linear digit-computation algorithms,
no such series is known for $\Kzero$. Every practical method expresses
it through the even zeta values $\zeta(2), \zeta(4), \ldots$ ---
infinitely many algebraically independent transcendentals, each needed
essentially to full precision --- and costs $O(d^2)$ up to logarithmic
factors for $d$ digits. This quadratic wall explains the brevity of the
record history (Table~\ref{tab:records}).

\begin{table}[ht]
\caption{Record computations of $\Kzero$.}
\label{tab:records}
\begin{tabular}{rrl}
\toprule
Year & Digits & Author / method \\
\midrule
1959--1960s & tens & Shanks, Wrench, and others \\
1997 & $7{,}350$ & Bailey, Borwein, Crandall \cite{bbc1997} \\
1998 & $110{,}000$ & Gourdon \cite{gourdon} \\
2016 & $1{,}000{,}000$ & Sim\'o \cite{simo2016}, PARI, ${\sim}12$ days, single core \\
2026 & $2{,}000{,}000$ & this work, 8\,h\,47\,m, one desktop machine \\
\bottomrule
\end{tabular}
\end{table}

\section{The series and its acceleration}

Taking logarithms of the product formula and expanding yields the
Bailey--Borwein--Crandall series \cite{bbc1997}
\begin{equation}
\label{eq:bbc}
\ln 2 \, \ln \Kzero
= \sum_{n=1}^{\infty} \frac{\zeta(2n) - 1}{n}\, h(n),
\qquad
h(n) = \sum_{j=1}^{2n-1} \frac{(-1)^{j+1}}{j}.
\end{equation}
Since $\zeta(2n) - 1 \sim 2^{-2n}$, the series
\eqref{eq:bbc} converges geometrically but slowly for high precision.
For a cut-off $N \ge 2$ the acceleration used here rewrites it as
\begin{equation}
\label{eq:accel}
\ln 2 \, \ln \Kzero
= -\sum_{k=2}^{N-1} \ln\frac{k-1}{k}\,\ln\frac{k+1}{k}
\;+\;
\sum_{n=1}^{\infty}
\frac{1}{n}\Bigl(\zeta(2n) - 1 - \sum_{k=2}^{N-1} k^{-2n}\Bigr) h(n),
\end{equation}
whose $n$-th term now has magnitude about $N^{-2n}$, so approximately
$P / (2\log_2 N)$ terms suffice for $P$ bits of precision. The
structure of this acceleration follows Bloemen's implementation
\cite{bloemen2022}.

\section{Implementation}

The program is a single C file over FLINT/Arb ball arithmetic
\cite{johansson2017} with OpenMP. Beyond the series
\eqref{eq:accel}, three ideas carry most of the speed.

\subsection*{Dropping precision}
A term of magnitude $N^{-2n}$ contributes at most
$P - 2n \log_2 N$ bits to the final sum, so its evaluation --- and the
power-table entries feeding it --- run far below full precision.
Almost all arithmetic in the program is low-precision.

\subsection*{A two-region split}
The bracketed quantity in \eqref{eq:accel} equals the tail
$\sum_{k \ge N} k^{-2n}$ exactly. For $n$ above roughly
$P / \bigl(2(\log_2 N + 3)\bigr)$ that tail is summed literally at low
precision, requiring no Bernoulli numbers at all; a rigorous integral
bound covers the truncated remainder. Below the split, where
$\zeta(2n)$ is genuinely near $1$ and needs full precision, FLINT's
reverse Bernoulli iterator reconstructs it. Work is divided into
cost-balanced blocks scheduled dynamically across all threads.

\subsection*{Rigor by construction}
Every intermediate quantity is an Arb ball --- an interval guaranteed
to contain the true value --- and every truncation is folded in as an
explicit error bound. The final decimal expansion is produced by a
proof-carrying rounding step that prints digits only when the ball
certifies a unique correctly rounded value. There is no estimated
error anywhere in the pipeline; a wrong digit is impossible, and
failure would be loud rather than silent.

\section{The record computation}

On August 27, 2026, the program computed $2{,}000{,}000$ digits after
the decimal point in a single run (Table~\ref{tab:run}).

\begin{table}[ht]
\caption{Parameters of the record run.}
\label{tab:run}
\begin{tabular}{ll}
\toprule
Hardware & Intel Core Ultra 9 275HX, 24 threads, 64\,GB RAM \\
Software & Fedora Linux 44, GCC 16.2, FLINT 3.4.0, stock GMP/MPFR \\
Working precision & $6{,}643{,}988$ bits \\
Wall time & 8\,h\,47\,m\,26\,s \quad (compute: $31{,}644$\,s) \\
Average CPU utilization & 931\% of 2400\% \\
Peak resident memory & 55.8\,GiB \\
\bottomrule
\end{tabular}
\end{table}

The run pressed against the machine's limits: the program's
conservative memory preflight refuses $2 \times 10^6$ digits on a
64\,GB machine (estimate 61.7\,GiB) and was overridden; the measured
peak of 55.8\,GiB left under 6\,GiB of headroom. Memory, not time, is
the binding constraint. The dominant allocation is the internal state
of the reverse Bernoulli iterators, which grows roughly quadratically
with the digit count; the measured local cost exponent, near
$d^{2.0}$ at moderate sizes, steepens to about $3.2$ between one and
two million digits as that state overwhelms the caches. A further
doubling would require well over 128\,GB of memory together with an
out-of-core redesign, and $10^9$ digits remains out of reach for every
known method: the exponent, not the constant factor, is the
obstruction.

\section{Verification}

Three independent legs support the result.
\begin{enumerate}
\item \emph{Interval certification.} The computation carries a
rigorous enclosure end to end; the printed digits are certified
correctly rounded by construction, not checked after the fact.
\item \emph{An independent backend.} The program ships a second,
structurally different backend (per-term \texttt{arb\_zeta\_ui} at
full precision) sharing only the outer series framing
\eqref{eq:accel}. The two agree byte for byte at every size tested.
\item \emph{Agreement with the previous record.} The first $10^6$
digits of the new result are identical, digit for digit, to Sim\'o's
independent 2016 computation \cite{simo2016} --- different program,
different algorithmic details, different hardware era --- as cached on
OEIS entry A002210.
\end{enumerate}

\section{Availability}

The C program, the two-million-digit result, the raw run telemetry,
the mathematical background document, and reference implementations of
the same accelerated series in seventeen further languages (each
verified byte-identical to the C program's output) are available under
the MIT license at
\url{https://github.com/reza-ghazi/khinchin-fast}
and are permanently archived at Zenodo
\cite{zenodo}. The computation is listed in the links of OEIS entry
A002210 \cite{oeis} as of August 31, 2026. A web version of this
article is available at
\url{https://reza-ghazi.github.io/khinchin-fast/}.

\section*{Acknowledgments}

The software and this write-up were developed with the assistance of
Claude (Anthropic), used as a programming and drafting tool under the
author's direction; all results are independently verifiable as
described in Section~5.

\begin{thebibliography}{9}

\bibitem{khinchin1935}
A.~Ya. Khinchin,
\emph{Metrische Kettenbruchprobleme},
Compositio Math. \textbf{1} (1935), 361--382.
See also \emph{Continued Fractions}, Dover, 1997.

\bibitem{bbc1997}
D.~H. Bailey, J.~M. Borwein, and R.~E. Crandall,
\emph{On the Khintchine constant},
Math. Comp. \textbf{66} (1997), no.~217, 417--431.

\bibitem{gourdon}
X.~Gourdon and P.~Sebah,
\emph{The Khintchine constant},
Numbers, constants and computation, 1998--2002.

\bibitem{simo2016}
C.~Sim\'o,
\emph{Computation of $10^6$ digits of Khintchine's constant}, 2016.
Announced via OEIS A002210;
\url{http://www.maia.ub.es/dsg/khinchin}.

\bibitem{oeis}
OEIS Foundation Inc.,
\emph{The On-Line Encyclopedia of Integer Sequences}, entry A002210,
\url{https://oeis.org/A002210}.

\bibitem{johansson2017}
F.~Johansson,
\emph{Arb: efficient arbitrary-precision midpoint-radius interval
arithmetic},
IEEE Trans. Comput. \textbf{66} (2017), no.~8, 1281--1292.
FLINT: \url{https://flintlib.org}.

\bibitem{bloemen2022}
R.~Bloemen,
\texttt{recmo/khinchin}, 2022,
\url{https://github.com/recmo/khinchin}.

\bibitem{zenodo}
R.~K~Ghazi,
\emph{khinchin-fast: Khinchin's constant to 2{,}000{,}000 certified
decimal digits}, Zenodo, 2026,
\url{https://doi.org/10.5281/zenodo.22134905}.

\end{thebibliography}

\end{document}
